% Beispielcodes für Abbildung 2.5 (Abschnitt 2.5, Maschinenlesbare Kennzeichnung).
% Erzeugt mit pst-barcode (BWIPP). pdflatex kann PSTricks nicht direkt, deshalb einmalig:
%   latex Kennzeichnung.tex
%   dvips Kennzeichnung.dvi
%   ps2pdf Kennzeichnung.ps Kennzeichnung.pdf
% width/height von \psbarcode sind in Zoll (2,2 cm = 0,866 in). pst-barcode zeichnet hier 20 pt unterhalb
% des Bezugspunkts, deshalb \rput(0,20pt). Jede pspicture wird eine eigene Seite, eingebunden per
% \includegraphics[page=n]. Prüfen: mgs -q -dNOPAUSE -dBATCH -sDEVICE=bbox Kennzeichnung.pdf
\documentclass[pstricks,border=0pt]{standalone}
\usepackage{pst-barcode}
\begin{document}
% Seite 1: Data Matrix des Platinen-Aufklebers (neues Format)
\begin{pspicture}(2.2,2.2)
	\rput[lb](0,20pt){\psbarcode{OEA101|10010310|00000001}{width=0.866 height=0.866}{datamatrix}}
\end{pspicture}
% Seite 2: QR-Code des Kistenetiketts, Fehlerkorrekturstufe M wie in der Anwendung
\begin{pspicture}(2.2,2.2)
	\rput[lb](0,20pt){\psbarcode{KISTE-000124}{width=0.866 height=0.866 eclevel=M}{qrcode}}
\end{pspicture}
% Seite 3: Code 128 des Moduletiketts (fiktive Beispielwerte)
\begin{pspicture}(5.4,1.5)
	\rput[lb](0,20pt){\psbarcode{10030001;L26010001}{width=2.126 height=0.591}{code128}}
\end{pspicture}
\end{document}
